\chapter{Design of the AI-Assisted Digitisation Pipeline}

The AI-assisted digitisation pipeline for historical inscriptions and manuscripts is built around a set of functional requirements, modular stage boundaries, and artefact-specific processing strategies that together translate the project's preservation goals into a concrete software architecture. This chapter presents the design decisions that underpin the Phase~1 system: the functional specifications and overall pipeline layout, the per-stage algorithmic choices and model selection rationale, the artefact-specific processing chains, a consolidated design decisions table, and the non-destructive processing policy that ensures original artefacts are never placed at risk. These design elements form a coherent blueprint that Chapter~4 then implements.

\section{Functional Requirements and Design Specifications}

The Phase~1 system must satisfy three core functional requirements derived from the project objectives:

\begin{enumerate}
    \item \textbf{Multi-artefact support:} Accept images from five artefact categories (stone inscriptions, palm leaf manuscripts, copper plates, paper manuscripts, and cave paintings), routing each through an artefact-specific processing chain.
    \item \textbf{Non-destructive processing:} Original images must never be overwritten. All outputs are written to separate directories, preserving the source artefact in its original form.
    \item \textbf{Scalable throughput:} The pipeline must process at least 10 images per minute on a single GPU to be practically useful for institutional-scale digitisation.
\end{enumerate}

The overall quality targets for Phase~1 are PSNR $\geq 30$ dB and SSIM $\geq 0.85$ across the test dataset. The modular design ensures that each stage can be upgraded or replaced independently without affecting the interface contracts of adjacent stages.

\section{Datasets Used for Development and Evaluation}

The pipeline is developed and evaluated using three primary datasets and four supplementary datasets drawn from Indian heritage archives and research institutions.

\subsection{Primary Datasets}

\begin{table}[htb]
\fontsize{10}{12}\selectfont
\centering
\caption{Primary datasets used for pipeline development and evaluation}
\label{tab:primary_datasets}
\begin{tabular}{|p{4.2cm}|p{5.0cm}|p{3.5cm}|}
\hline
\textbf{Dataset} & \textbf{Contents} & \textbf{Source} \\
\hline
Ancient Tamil Stone Inscriptions & Tamil stone inscription photographs with LiDAR scans and 3D models & Kaggle (Athira Ishanmugam) \\
\hline
Tamil Handwritten Palm Leaf (THPLMD) \cite{Dhanya2020} & 262 deteriorated Tamil palm leaf manuscript samples with binarised ground truth annotations & ScienceDirect / PMC \\
\hline
Tamil Handwritten Character Corpus & Tamil handwritten characters spanning multiple centuries and scribal styles & data.mendeley.com \\
\hline
\end{tabular}
\end{table}

The THPLMD dataset provides the most rigorous evaluation baseline because it includes expert-validated binarised ground truth, enabling direct comparison of pipeline binarisation output against a known-good reference. The Kaggle Tamil Stone Inscriptions dataset includes 3D point cloud models alongside photographs, supporting future Phase 3 work on depth-enhanced enhancement.

\subsection{Secondary and Supplementary Datasets}

Four additional datasets inform Phase~2 planning and script extension work:

\begin{table}[htb]
\fontsize{10}{12}\selectfont
\centering
\caption{Secondary datasets for extended script and layout support}
\label{tab:secondary_datasets}
\begin{tabular}{|p{4.2cm}|p{5.0cm}|p{3.5cm}|}
\hline
\textbf{Dataset} & \textbf{Contents} & \textbf{Source} \\
\hline
Indiscapes (IIIT Hyderabad) \cite{Mehrotra2022} & Layout annotations for historical Indic manuscripts & ihdia.iiit.ac.in \\
\hline
Sanskrit OCR Dataset & Classical Sanskrit document images & ihdia.iiit.ac.in \\
\hline
Brahmi Character Dataset \cite{Aitha2023} & 1032 Ashokan Brahmi characters, 258 classes & arXiv:2501.01981 \\
\hline
Kannada Inscriptions & Palm leaf manuscripts and stone inscriptions from the Hampi region & Kuvempu Institute of Kannada Studies \\
\hline
\end{tabular}
\end{table}

The Brahmi Character Dataset is the planned training resource for a custom Phase~2 Brahmi character recognition model. The Indiscapes dataset provides layout annotation methodology that can extend future page segmentation capabilities for multi-column Indic manuscript layouts.

\subsection{Institutional Archives}

Institutional archives form the long-term source of artefact images for production use. Key archives include the Archaeological Survey of India (asi.nic.in), the Digital Library of India (dli.ernet.in), the eGangotri portal (egangotri.org), and the IIIT Hyderabad IHDIA collection (ihdia.iiit.ac.in). All archived images are in standard JPEG, PNG, or TIFF formats compatible with the pipeline's input specification.

\section{Overall Pipeline Architecture}

The Phase~1 pipeline is structured as three sequential stages. Table~\ref{tab:pipeline_overview} summarises each stage, its processing operation, and its output format.

\begin{table}[htb]
\fontsize{10}{12}\selectfont
\centering
\caption{Phase~1 pipeline stage overview: input, processing, and output}
\label{tab:pipeline_overview}
\begin{tabular}{|c|p{3.2cm}|p{4.2cm}|p{3.2cm}|}
\hline
\textbf{Stage} & \textbf{Name} & \textbf{Processing} & \textbf{Output} \\
\hline
1 & Preprocessing & CLAHE, white balance, EXIF orientation, border crop & Normalised JPEG (quality=95) \\
\hline
2 & AI Enhancement & Real-ESRGAN, denoising, DStretch, sharpening & Enhanced JPEG \\
\hline
3 & Binarisation & Sauvola adaptive thresholding, morphological closing & Binary PNG \\
\hline
\end{tabular}
\end{table}

Figure~\ref{fig:pipeline_architecture} provides a visual overview of the pipeline flow.

\begin{figure}[H]
\centering
\begin{tikzpicture}[
    box/.style={draw=black!65, fill=blue!12, rounded corners=4pt,
                minimum width=3.2cm, minimum height=1.0cm, align=center,
                font=\small\sffamily},
    ph2box/.style={draw=gray!60, fill=gray!8, rounded corners=4pt,
                minimum width=3.6cm, minimum height=1.0cm, align=center,
                font=\small\sffamily, dashed},
    arr/.style={-Stealth, thick, draw=black!55}
  ]
  \node[box]    (s1) at (0,0)   {\textbf{Stage 1}\\Preprocessing};
  \node[box]    (s2) at (4.0,0) {\textbf{Stage 2}\\AI Enhancement};
  \node[box]    (s3) at (8.0,0) {\textbf{Stage 3}\\Binarisation};
  \node[ph2box] (s4) at (12.2,0){\textbf{Stage 4+}\\(Phase~2)};
  \node[font=\small\itshape] at (-1.8,0) {Raw scan};
  \draw[arr] (s1) -- (s2);
  \draw[arr] (s2) -- (s3);
  \draw[arr] (s3) -- (s4);
\end{tikzpicture}
\caption{Phase~1 pipeline (three solid stages) feeding into deferred Phase~2 stages. All solid boxes are fully implemented; the dashed box represents downstream stages planned for Phase~2.}
\label{fig:pipeline_architecture}
\end{figure}

Stages 1--3 are fully implemented in Phase~1. Stage~4 and all subsequent stages are deferred to Phase~2 and beyond, per mentor guidance to first deliver a robust, well-evaluated image enhancement and binarisation system.

\subsection{Design Principles}

Three principles govern every architectural decision in the system.

\textbf{Non-destructive preservation:} The \texttt{data/raw/} directory is read-only at all times. Every processing stage writes exclusively to its own output directory (\texttt{data/enhanced/}, \texttt{data/binarised/}, etc.). No original image file is ever modified or overwritten.

\textbf{Modular, artefact-aware routing:} Rather than applying a single fixed chain to every input, the pipeline selects enhancement parameters appropriate for each artefact type. This routing is implemented as a configuration look-up keyed on the artefact type string provided by the user at submission time, making it easy to add new artefact categories in future without modifying existing chains.

\textbf{Quality-transparent outputs:} Every pipeline output carries associated quality measures. Image quality is reported as PSNR, SSIM, CNR, and sharpness after Stage~2 enhancement. This allows downstream users and automated checks to identify images requiring additional manual attention before Phase~2 processing.

\section{Stage-by-Stage Design}

\subsection{Stage 1 --- Preprocessing Design}

The preprocessing stage normalises raw scanned images to a consistent quality baseline before any AI processing occurs. The design specifies four operations applied in sequence:

\begin{itemize}
    \item \textbf{EXIF orientation correction:} Camera images carry an EXIF orientation tag indicating how the image was rotated at capture. The design requires this tag to be baked into pixel data before saving so that all downstream stages receive a geometrically upright image regardless of capture orientation. This is performed using \texttt{PIL.ImageOps.exif\_transpose}.
    \item \textbf{Brightness normalisation:} Contrast Limited Adaptive Histogram Equalisation (CLAHE) with \texttt{clipLimit=2.0} and \texttt{tileGridSize=(8, 8)} is applied to the luminance channel in LAB colour space. Applying CLAHE only to luminance avoids introducing colour artefacts while correcting spatially uneven illumination from field photography.
    \item \textbf{White balance:} A grey-world assumption correction scales each colour channel's gain by the ratio of the global mean intensity to the per-channel mean. This removes systematic colour casts from scanner lamps or sunlight without requiring a reference white patch.
    \item \textbf{Border removal:} Rows and columns where the mean pixel intensity falls below a threshold of 10 (near-black) are identified as scanner margins and removed. This prevents scanner beds and dark borders from interfering with histogram-based enhancement in Stage 2.
\end{itemize}

The output format is JPEG at quality=95. TIFF was evaluated but rejected due to metadata write failures with the PIL/libtiff combination on Windows hosts. Quality=95 JPEG is indistinguishable from lossless for all subsequent pipeline operations.

\subsection{Stage 2 --- AI Enhancement Design}

The enhancement stage is the core value-add of the system: it transforms previously illegible scans into visually recoverable images. Three components are composed into artefact-specific chains.

\textbf{Denoising:} Non-local means (NLM) denoising is selected over Gaussian blur or median filtering because NLM compares patches across the whole image to average similar regions, preserving text edge structure while suppressing noise. The denoising strength is a per-artefact parameter informed by ECE noise characterisation: 10 for stone inscriptions (Gaussian and salt-and-pepper noise at moderate levels), 15 for palm leaf manuscripts (heavier Gaussian noise from yellowed substrate), and 8 for copper plates (lower strength to preserve metallic texture detail).

\textbf{Super-resolution --- Real-ESRGAN:} Bicubic upscaling cannot recover high-frequency character detail from degraded scans. Real-ESRGAN \cite{Wang2021}, trained on pairs of high-quality photographs and their realistically degraded counterparts, learns the inverse of real-world degradation functions and produces textured, character-aware outputs. The RRDB architecture (23 residual dense blocks, 64 feature maps) captures both global structure and fine character strokes. The design uses \texttt{RealESRGAN\_x4plus.pth} for most inscription types and \texttt{RealESRGAN\_x4plus\_anime\_6B.pth} for line-art carving styles. The output scale is set to 2 (using the 4$\times$ model internally but downsampling by a factor of 2), which avoids hallucinating detail that is not present in the original scan --- a critical requirement when fidelity to the historical source is paramount.

\textbf{DStretch:} Decorrelation stretch, developed by Jon Harman \cite{Harman2008} for rock art analysis, reveals colour information imperceptible to human vision by removing inter-channel correlation. The eigendecomposition of the pixel colour covariance matrix whitens the colour space, amplifying subtle colour differences. This is particularly powerful for stone inscriptions and cave paintings where the visual distinction between inscription and substrate is too faint for standard enhancement. LAB colour space is used for stone; YBK for cave and rock art.

\textbf{Sharpening:} Unsharp masking with amount=1.5 is applied after super-resolution to crisp character edge definition, compensating for any residual blur from the denoising or upscaling steps.

The per-artefact processing chains are designed as follows:
\begin{itemize}
    \item Stone inscriptions: denoise $\rightarrow$ Real-ESRGAN $\rightarrow$ sharpen
    \item Cave / rock paintings: denoise $\rightarrow$ DStretch (YBK) $\rightarrow$ sharpen
    \item Palm leaf manuscripts: white balance $\rightarrow$ denoise $\rightarrow$ Real-ESRGAN
    \item Copper plate inscriptions: HDR normalisation $\rightarrow$ denoise $\rightarrow$ Real-ESRGAN
    \item Paper manuscripts: denoise $\rightarrow$ Real-ESRGAN $\rightarrow$ sharpen
\end{itemize}

\subsection{Stage 3 --- Binarisation Design}

Binarisation converts the enhanced colour image to a clean binary image (black text on white background), which is the Phase~1 endpoint and the input for Phase~2 character recognition. Three methods are designed, with a preferred default:

\textbf{Primary method --- Sauvola adaptive thresholding \cite{Dhanya2020}:} Sauvola's local threshold adapts to local pixel statistics within a 25$\times$25 pixel window, using the local mean and standard deviation to compute a pixel-specific threshold. This handles spatially uneven backgrounds (stone surface texture, aged palm leaf fibre) far better than global methods. Sauvola is selected as the default for all five artefact types.

\textbf{Secondary --- Otsu global thresholding:} Otsu's method maximises between-class variance to select a single global threshold from the histogram. It is appropriate for clean paper manuscripts where the background is uniform and the bimodal histogram assumption holds.

\textbf{Fallback --- OpenCV adaptive mean thresholding:} Available for images of mixed quality where Sauvola produces excessive noise due to extreme local intensity variation.

Post-binarisation, two refinement steps are applied: a 2$\times$2 morphological closing kernel (\texttt{cv2.MORPH\_CLOSE}) reconnects broken character strokes caused by surface erosion on stone inscriptions; and a minimum blob size filter (50 pixels) removes dust particles and surface damage artefacts that appear as disconnected components too small to represent character strokes. The output format is PNG (lossless), which introduces no compression artefacts on binary data --- an important distinction from JPEG, which would degrade the clean text/background boundary.

\section{Key Design Decisions and Rationale}

Table~\ref{tab:design_decisions} consolidates the major design choices made during the architecture phase and the reasoning behind each selection.

\begin{table}[htb]
\fontsize{10}{12}\selectfont
\centering
\caption{Key design decisions and rationale}
\label{tab:design_decisions}
\begin{tabular}{|p{4.5cm}|p{8.2cm}|}
\hline
\textbf{Design Decision} & \textbf{Rationale} \\
\hline
Real-ESRGAN over bicubic upscaling & Trained on real-world degraded photographs; recovers character stroke texture that polynomial interpolation cannot reconstruct \\
\hline
DStretch for cave art and stone & Eigendecomposition-based colour decorrelation reveals features invisible to the human eye; originally developed specifically for rock art analysis \\
\hline
Sauvola as default binarisation & Adapts to spatially uneven backgrounds (stone texture, palm leaf fibre); Otsu's global threshold fails on non-uniform illumination \\
\hline
JPEG quality=95 for preprocessing & Avoids PIL/libtiff compatibility failures on Windows; quality=95 is indistinguishable from lossless for all downstream stages \\
\hline
\texttt{outscale=2} for Real-ESRGAN & Uses the 4$\times$ model internally but outputs at 2$\times$ to avoid hallucinating detail not present in the original inscription \\
\hline
Phase~1 endpoint at binarisation & Mentor guidance: deliver robust, validated image enhancement and binarisation first; ensures Phase~2 receives the highest possible input quality \\
\hline
\end{tabular}
\end{table}

\section{Non-Destructive Processing Policy}

The non-destructive processing policy is a mandatory system constraint adopted from archival preservation standards. It applies to all pipeline stages without exception and encompasses five rules:

\begin{enumerate}
    \item \textbf{Read-only originals:} The \texttt{data/raw/} directory is never written to by any pipeline process. All enhanced and binarised outputs go to designated sibling directories (\texttt{data/enhanced/}, \texttt{data/binarised/}).
    \item \textbf{3-2-1 backup rule:} Three copies of original data are maintained on two different storage media, with one copy held offsite or in cloud storage.
    \item \textbf{EXIF metadata preservation:} Orientation and capture metadata extracted at Stage~1 are re-embedded into all output images throughout the pipeline.
    \item \textbf{Full processing log:} A processing log is maintained recording the pipeline stage, timestamp, software version, and model weights version for each operation, providing a complete and reproducible audit trail.
    \item \textbf{Model version control:} The exact Real-ESRGAN weights file used for each enhancement is recorded in the processing log, enabling future re-processing with updated models while preserving provenance of every existing output.
\end{enumerate}

This policy ensures that no irreplaceable historical source is damaged or lost through the digitisation process, and that every digital record produced is fully reproducible and attributable to a specific software configuration.

\section{Summary}

The design of the Phase~1 digitisation pipeline is built on functional requirements that address the practical needs of heritage institutions: multi-artefact processing, non-destructive handling of originals, and quality-transparent outputs. The three-stage architecture provides clear interface boundaries between preprocessing, AI enhancement, and binarisation, enabling modular improvement of any single stage. The key design decisions --- Real-ESRGAN for super-resolution, DStretch for colour-faint inscriptions, and Sauvola for adaptive binarisation --- are each grounded in the specific degradation characteristics of South Asian heritage artefacts. The non-destructive processing policy and objective quality metrics ensure that the system produces validated, reproducible outputs suitable for institutional archival workflows and ready for Phase~2 processing. Chapter~4 now details how each of these design decisions is realised in working code.
